\section{Use Case Diagram e narrative}

\subsection{Use Case Diagram}
\begin{figure}[H]
    \centering
    % \includegraphics[width=0.85\textwidth]{use-case-diagram.png}
    \fbox{\parbox{0.85\textwidth}{\centering\vspace{1em}
        \TODO{Inserire \emph{Use Case Diagram} UML disegnato con draw.io o
        Excalidraw. Attori: ``Utente anonimo'' e ``Utente registrato'';
        casi d'uso: Cerca brano, Esplora catalogo, Legge testo
        sincronizzato, Registrazione, Login, Pubblica anonimo (PoW),
        Pubblica autenticato, Modifica brano, Cancella brano.}
        \vspace{1em}}}
    \caption{Use Case Diagram di Stopify.}
    \label{fig:usecase}
\end{figure}

\subsection{Narrative dei casi d'uso}

% =====================================================================
\paragraph{UC1 --- Ricerca di un brano}\hfill\\
\textbf{Scope:} catalogo lyrics.\\
\textbf{Level:} user-goal.\\
\textbf{Intention in context:} l'utente vuole leggere il testo di un brano
che gli è venuto in mente.\\
\textbf{Primary actor:} Utente anonimo.\\
\textbf{Support actor:} Backend Flask.

\subparagraph{Stakeholder's interest:}
\begin{itemize}
    \item \textbf{Utente:} trovare velocemente il testo di un brano.
    \item \textbf{Sistema:} servire risultati rilevanti con bassa latenza.
\end{itemize}

\textbf{Precondition:} il backend è raggiungibile.\\
\textbf{Minimum guarantees:} nessuna modifica dello stato del sistema.\\
\textbf{Success guarantees:} l'utente vede il testo del brano richiesto.\\
\textbf{Trigger:} l'utente apre la schermata ``Cerca''.

\subparagraph{Main success scenario:}
\begin{enumerate}
    \item L'utente digita una stringa nel campo di ricerca.
    \item L'app invia \endpoint{GET}{/api/search?q=<stringa>}.
    \item Il backend restituisce l'elenco dei brani che matchano titolo o
    artista.
    \item L'utente seleziona un brano.
    \item L'app invia \endpoint{GET}{/songs/<id>} e mostra il testo.
\end{enumerate}

\subparagraph{Extensions:}
\begin{itemize}
    \item \textbf{3a:} nessun risultato $\Rightarrow$ l'app mostra empty state
    ``Nessun risultato''.
\end{itemize}

% =====================================================================
\paragraph{UC2 --- Registrazione di un nuovo account}\hfill\\
\textbf{Scope:} autenticazione.\\
\textbf{Level:} user-goal.\\
\textbf{Intention in context:} l'utente vuole creare un account per
contribuire senza dover risolvere la PoW ad ogni pubblicazione.\\
\textbf{Primary actor:} Utente anonimo.

\textbf{Precondition:} l'utente non è autenticato.\\
\textbf{Success guarantees:} è stato creato un nuovo record \code{User} e
l'utente è loggato.

\subparagraph{Main success scenario:}
\begin{enumerate}
    \item L'utente apre la schermata ``Crea un account''.
    \item Inserisce username, email e password.
    \item L'app valida i campi lato client (regex, lunghezza minima).
    \item L'app invia \endpoint{POST}{/auth/register} con il payload.
    \item Il backend valida nuovamente, hasha la password (PBKDF2-SHA256),
    crea il record e restituisce \emph{access} e \emph{refresh token}.
    \item L'app salva i token nello storage sicuro e mostra la home da
    autenticato.
\end{enumerate}

\subparagraph{Extensions:}
\begin{itemize}
    \item \textbf{4a:} username o email già presenti $\Rightarrow$ il backend
    risponde 409 e l'app mostra il messaggio di errore inline.
    \item \textbf{5a:} password troppo debole $\Rightarrow$ 422 con
    \code{WeakPasswordError}.
\end{itemize}

% =====================================================================
\paragraph{UC3 --- Pubblicazione anonima con Proof of Work}\hfill\\
\textbf{Scope:} contenuti, sicurezza.\\
\textbf{Level:} user-goal.\\
\textbf{Intention in context:} l'utente vuole condividere il testo di un
brano senza voler creare un account.\\
\textbf{Primary actor:} Utente anonimo.

\textbf{Precondition:} l'utente conosce titolo, artista e testo.\\
\textbf{Minimum guarantees:} nessun brano viene pubblicato se la PoW non
viene superata.\\
\textbf{Success guarantees:} un nuovo record \code{Song} è presente nel DB
con \code{user\_id=NULL}.

\subparagraph{Main success scenario:}
\begin{enumerate}
    \item L'utente apre la schermata ``Pubblica'' e compila il form.
    \item L'app invia \endpoint{POST}{/api/request-challenge} e riceve
    \code{\{prefix, target\}}.
    \item L'app risolve la PoW cercando un \code{nonce} tale che
    SHA-256(\code{prefix+nonce}) $\leq$ \code{target}.
    \item L'app invia \endpoint{POST}{/api/publish} con il body del brano e
    l'header \code{X-Publish-Token: prefix:nonce}.
    \item Il backend verifica la PoW, marca atomicamente il token come
    consumato, crea il brano e risponde 201 con il record creato.
\end{enumerate}

\subparagraph{Extensions:}
\begin{itemize}
    \item \textbf{4a:} il token PoW è già stato consumato (race condition)
    $\Rightarrow$ 409 \code{Conflict}.
    \item \textbf{4b:} il \code{plainLyrics} eccede la lunghezza massima
    $\Rightarrow$ 422 con dettaglio della violazione.
\end{itemize}

% =====================================================================
\paragraph{UC4 --- Pubblicazione autenticata}\hfill\\
\textbf{Scope:} contenuti.\\
\textbf{Level:} user-goal.\\
\textbf{Primary actor:} Utente registrato.

\textbf{Precondition:} l'utente è autenticato (\emph{access token} valido).\\
\textbf{Success guarantees:} il nuovo brano viene creato con
\code{user\_id} valorizzato.

\subparagraph{Main success scenario:}
\begin{enumerate}
    \item L'utente apre la schermata ``Pubblica''.
    \item L'app, vedendo il contesto autenticato, salta la richiesta di
    challenge.
    \item L'app invia \endpoint{POST}{/api/publish} con header
    \code{Authorization: Bearer <jwt>}.
    \item Il backend identifica l'utente dal token e crea il brano con
    \code{user\_id} valorizzato.
\end{enumerate}

% =====================================================================
\paragraph{UC5 --- Modifica di un brano pubblicato}\hfill\\
\textbf{Scope:} contenuti, autorizzazione.\\
\textbf{Level:} user-goal.\\
\textbf{Primary actor:} Utente registrato.

\textbf{Precondition:} l'utente è autenticato e proprietario del brano.\\
\textbf{Minimum guarantees:} un brano non può essere modificato da un
utente diverso dal proprio autore.

\subparagraph{Main success scenario:}
\begin{enumerate}
    \item L'utente apre ``I miei testi''.
    \item L'app chiama \endpoint{GET}{/api/me/songs} e mostra l'elenco.
    \item L'utente tocca l'icona di modifica su un brano.
    \item L'app naviga al form precompilato.
    \item Dopo le modifiche, invia \endpoint{PUT}{/api/songs/<id>}.
    \item Il backend verifica che \code{song.user\_id == current\_user\_id},
    aggiorna i campi e risponde 200.
\end{enumerate}

\subparagraph{Extensions:}
\begin{itemize}
    \item \textbf{6a:} \code{song.user\_id} appartiene a un altro utente
    $\Rightarrow$ 403 \code{Forbidden}.
    \item \textbf{6b:} il brano è anonimo (\code{user\_id=NULL})
    $\Rightarrow$ 403 (non modificabile).
\end{itemize}

% =====================================================================
\paragraph{UC6 --- Cancellazione di un brano}\hfill\\
\textbf{Scope:} contenuti, autorizzazione.\\
\textbf{Primary actor:} Utente registrato.

\textbf{Success guarantees:} il record \code{Song} viene rimosso dal DB e
non è più reperibile via API.

\subparagraph{Main success scenario:}
\begin{enumerate}
    \item L'utente apre ``I miei testi''.
    \item Tocca l'icona di cancellazione.
    \item L'app mostra un alert di conferma.
    \item Su conferma, invia \endpoint{DELETE}{/api/songs/<id>}.
    \item Il backend verifica la proprietà e cancella; risponde 200.
\end{enumerate}

\subparagraph{Extensions:}
\begin{itemize}
    \item \textbf{4a:} l'utente annulla l'alert $\Rightarrow$ nessuna azione.
    \item \textbf{5a:} brano non di proprietà $\Rightarrow$ 403.
\end{itemize}
